% =====================================================================
%  A State-Resolved Calculus of Metabolic Probes:
%  Diagnosis as Localization, Therapy as Gap-Closure, and a
%  Domain-Specific Language over a First-Principles Pathway Runtime
%
%  Self-contained. No companion manuscript is presupposed. Every object
%  is a dynamical system on a finite-dimensional state manifold or a
%  construction on one; every claim is stated with its hypotheses.
% =====================================================================
\documentclass[11pt,a4paper]{article}

\usepackage[utf8]{inputenc}
\usepackage[T1]{fontenc}
\usepackage{lmodern}
\usepackage[a4paper,margin=2.5cm]{geometry}
\usepackage{amsmath,amssymb,amsthm,mathtools}
\usepackage{bm}
\usepackage{mathrsfs}
\usepackage{booktabs}
\usepackage{enumitem}
\usepackage{microtype}
\usepackage[colorlinks=true,linkcolor=blue!45!black,
            citecolor=blue!45!black,urlcolor=blue!45!black]{hyperref}
\usepackage{cleveref}

% ---- Theorem environments ----
\newtheoremstyle{thmstyle}{\topsep}{\topsep}{\itshape}{}{\bfseries}{.}{.5em}{}
\theoremstyle{thmstyle}
\newtheorem{theorem}{Theorem}[section]
\newtheorem{lemma}[theorem]{Lemma}
\newtheorem{proposition}[theorem]{Proposition}
\newtheorem{corollary}[theorem]{Corollary}
\theoremstyle{definition}
\newtheorem{definition}[theorem]{Definition}
\newtheorem{assumption}{Assumption}
\newtheorem{principle}[theorem]{Principle}
\newtheorem{problem}[theorem]{Problem}
\theoremstyle{remark}
\newtheorem{remark}[theorem]{Remark}

\crefname{assumption}{Assumption}{Assumptions}
\Crefname{assumption}{Assumption}{Assumptions}

% ---- Shorthand ----
\newcommand{\R}{\mathbb{R}}
\newcommand{\N}{\mathbb{N}}
\newcommand{\Rp}{\R_{\ge 0}}
\newcommand{\Man}{\mathcal{M}}          % state manifold
\newcommand{\st}{x}                     % state
\newcommand{\coh}{x^\star}              % coherent state
\newcommand{\fld}{F}                    % vector field / dynamics
\newcommand{\probe}{\pi}                % probe program
\newcommand{\trace}{\tau}               % trace / readout
\newcommand{\Rd}{\mathcal{R}}           % readout map
\newcommand{\loc}{\Lambda}              % localization map
\newcommand{\gap}{\mathsf{g}}           % gap
\newcommand{\flo}{\beta}                % floor
\newcommand{\basis}{\mathcal{B}}        % intervention basis
\newcommand{\intv}{u}                   % intervention
\newcommand{\dsl}{\mathcal{L}}          % the DSL
\newcommand{\abs}[1]{\lvert#1\rvert}
\newcommand{\norm}[1]{\lVert#1\rVert}
\newcommand{\set}[1]{\{#1\}}
\newcommand{\eps}{\varepsilon}
\DeclareMathOperator*{\argmin}{arg\,min}
\DeclareMathOperator{\supp}{supp}
\DeclareMathOperator{\dist}{dist}

\title{\textbf{A State-Resolved Calculus of Metabolic Probes:\\[3pt]
Diagnosis as Localization, Therapy as Gap-Closure, and a\\
Language over a First-Principles Pathway Runtime}}

\author{}
\date{}

\begin{document}
\maketitle

% =====================================================================
\begin{abstract}
\noindent
We give a self-contained mathematical description of a diagnostic and
therapeutic method that dispenses with the two classical primitives of
clinical practice---the patient's report and the named disease---and
replaces them with a single primitive: a \emph{characterized
perturbation} of metabolism whose fate is computed from a state-resolved
forward model. The method rests on three observations, each made precise
below. First, a finite observer of a biochemical system cannot resolve
its state to a point; there is a strictly positive floor $\flo>0$ below
which distinctions cannot be made, so ``the state'' is always a region of
diameter $\ge\flo$, and the goal of measurement is not to name the state
but to shrink the region to the floor. Second, an ingested compound of
\emph{known} composition is a controlled input to the metabolic
dynamical system, and its fate---the time-course of products it
generates---is a function of the patient's state; if the forward model
that predicts this fate is \emph{state-resolved} (valid across the state
manifold, not merely at a healthy reference), then the measured fate
\emph{locates} the patient on the manifold. Third, because each product
of a metabolic reaction is itself a reactive species, a designed probe
can be made to \emph{cascade}: its first product is the substrate of the
next reaction, and so on, so that a single ingestion installs a
self-propagating chain whose later products \emph{exist only if} the
earlier trajectory occurred as predicted. The measured product
distribution is therefore not a snapshot but a \emph{certificate} of a
trajectory: a false account of the state cannot forge it, because it
would have to produce products that never formed.

From these we develop: (i) a \textbf{Floor Theorem} giving the positive
resolution limit and its consequence that ``solved'' is a region; (ii) a
\textbf{Localization Theorem} stating that a probe is diagnostically
sufficient iff the map (state~$\to$~trajectory~$\to$~measured product
distribution) is injective down to the floor, so that reading the
distribution pins the state to a floor-width cell; (iii) a
\textbf{Certificate Theorem} formalizing why a cascading probe's
terminal products constitute a proof of the path that generated them,
under which no state-account inconsistent with the observed products
survives; (iv) a \textbf{Gap-Closure} formulation of therapy in which the
disease is never named: the therapeutic target is the vector from the
measured state to the nearest coherent state, and treatment is the
sparsest combination of characterized interventions---drawn from a basis
of compounds with known state-resolved effects---whose local composition
drives that vector to the floor, solved as an $\ell_1$-regularized program
and executed as a contraction (measure, apply, re-measure) whose fixed
point is coherence; and (v) a proof that \emph{diagnosis and therapy are
one operation}, since every characterized perturbation simultaneously
extracts information (its fate is a measurement) and displaces the state
(its action is an intervention). We prove a \textbf{Stopping Theorem}: the
correct halting criterion is not a magnitude threshold but a
\emph{rigidity} condition---one halts when the accumulated products
over-determine the state, i.e.\ when every account other than the true one
fails to reproduce some observed product, and the surviving accounts differ
only within the floor-width cell.

Finally we describe the computational surface: because the forward model
is a runnable dynamical system, the entire method is a
\textbf{domain-specific language} over a \emph{first-principles} pathway
runtime. A user specifies a probe (a program of reactions) and an intent
(a target region); the runtime executes the probe against a model of the
patient's pathways and returns the predicted trace, the implied
localization, the gap, and the minimal intervention. The essential
requirement, which we state as a theorem rather than assume, is that the
runtime be \emph{derived} rather than \emph{fitted}: a parameter-fitted
model is valid only near its training regime and becomes unreliable
exactly where free exploration drives it, whereas a model whose rate laws
are fixed by physical law admits extrapolation across the manifold and so
supports risk-free \emph{in silico} experimentation---the user searches
the intervention space against the model, and only the chosen intervention
is applied to the patient. We give the language's core judgments, the
soundness conditions the runtime must satisfy, the complexity of each
operation, and a constructive validation protocol. The result is a
deterministic, disease-free, report-free calculus of medicine in which
the patient's state is read from a self-certifying metabolic program, the
therapeutic target is a measured gap, and the physician's role is
relocated from precedent-retrieval to judgment over a derived state.
\end{abstract}

\vspace{4pt}
\noindent\textbf{Keywords:} state-resolved forward model; metabolic probe;
reaction cascade; localization; resolution floor; certificate of trajectory;
gap-closure; sparse intervention; contraction to a fixed point; coherent
state; domain-specific language; first-principles runtime; \emph{in silico}
experimentation.

\tableofcontents

% =====================================================================
\section{Setting and Primitives}
\label{sec:setting}
% =====================================================================

We fix a mathematical setting in which metabolism is a controlled
dynamical system, an ingested compound is an input, and a measurement is
a readout of the induced trajectory. Everything in the paper is a
construction on the objects of this section.

\begin{definition}[Pathway state system]
\label{def:system}
A \emph{pathway state system} is a tuple
$\Sigma=(\Man,\ \fld,\ \mathcal{U},\ \Rd)$ where:
\begin{itemize}[leftmargin=1.4em,itemsep=2pt]
\item $\Man\subseteq\Rp^{n}$ is the \emph{state manifold}: a point
$\st\in\Man$ is a full description of the relevant biochemistry
(concentrations of metabolites, enzyme activities, cofactor levels,
transporter states) at an instant. The dimension $n$ is finite but may be
large.
\item $\fld\colon\Man\times\mathcal{U}\to\R^{n}$ is the \emph{dynamics}:
$\dot\st=\fld(\st,\intv)$ with $\intv$ a control input. $\fld$ is the
biochemical rate law---mass action, Michaelis--Menten, or any kinetic law
whose form is fixed by the chemistry and whose parameters are, by
hypothesis (\cref{sec:runtime}), determined by physical constants rather
than fitted.
\item $\mathcal{U}$ is the \emph{input space} of admissible perturbations
(ingested compounds and their dosing profiles).
\item $\Rd\colon\Man\to\mathcal{Y}$ is the \emph{readout map} into an
observation space $\mathcal{Y}$: what an instrument reports of a state.
The canonical readout is a mass or spectral distribution over species,
$\Rd(\st)=\bigl(\text{abundance of species } s \text{ in state } \st\bigr)_{s}$.
\end{itemize}
\end{definition}

\begin{definition}[Probe]
\label{def:probe}
A \emph{probe} $\probe$ is an element of $\mathcal{U}$ of \emph{known
composition}: a compound (or timed sequence of compounds) whose molecular
identity and quantity are exactly specified. The distinguishing property
of a probe, as against an arbitrary ingested substance, is that its
identity is known, so its fate under $\fld$ is \emph{predictable} given
the state $\st$.
\end{definition}

\begin{definition}[Fate; trace]
\label{def:fate}
Given a state $\st_0\in\Man$ and a probe $\probe$, the \emph{fate} of
$\probe$ is the trajectory $\st(t;\st_0,\probe)$ obtained by integrating
$\dot\st=\fld(\st,\probe)$ from $\st_0$. The \emph{trace} of the probe is
the time-indexed readout
\[
\trace(\st_0,\probe)\;=\;\bigl(\Rd(\st(t;\st_0,\probe))\bigr)_{t\ge 0},
\]
the observable record of the fate. In practice one reads the trace at one
or several times; a single spectral measurement returns the product
distribution $\Rd(\st(t^\ast))$ at the read time $t^\ast$.
\end{definition}

\begin{remark}[A probe is a meal of known composition]
\label{rem:meal}
Nothing in \cref{def:probe} requires the probe to be exotic or inert. Any
ingested substance is a perturbation of $\Sigma$; a meal is such a
perturbation of \emph{unknown} composition, hence of unpredictable fate. A
probe is a perturbation of the same physical class whose composition is
known, and it is exactly this knowledge---not any special inertness---that
makes its fate a usable measurement. The perturbation is not a hazard to be
avoided but the instrument itself.
\end{remark}

We adopt two structural assumptions and carry them throughout.

\begin{assumption}[Finite resolution]
\label{ass:finite}
Any physical observer of $\Sigma$ has finite resolving power: the readout
$\Rd$ distinguishes states only up to a positive tolerance. Formally there
is a $\delta>0$ such that states within distance $\delta$ produce
observationally indistinguishable readouts,
$\dist(\st,\st')<\delta \Rightarrow \norm{\Rd(\st)-\Rd(\st')}$ below the
instrument's discrimination threshold.
\end{assumption}

\begin{assumption}[Non-exhaustibility]
\label{ass:nonexhaust}
No finite measurement exhausts the state: for any finite battery of
readouts there remains a distinction the system could exhibit that the
battery does not resolve. Equivalently, the state manifold admits, at every
scale, a finer sub-state the current resolution collapses. This is an order
condition on refinement, not a claim of infinite dimension.
\end{assumption}

% =====================================================================
\section{The Resolution Floor}
\label{sec:floor}
% =====================================================================

The first quantitative fact is that resolution has a positive floor, and
that this floor is not a nuisance to be engineered away but the width of
the region within which a state counts as determined.

\begin{definition}[Separation and residue]
\label{def:residue}
For a state $\st$, its \emph{separation} $\res(\st)$ is the least
observational effort---measured as the discrimination distance in
$\mathcal{Y}$ realized over admissible probes---required to tell $\st$
apart from every other state of $\Man$. The \emph{residue} of $\st$ is the
irreducible part of this separation that no admissible measurement removes.
\end{definition}

\begin{theorem}[Floor]
\label{thm:floor}
Under \cref{ass:finite,ass:nonexhaust} there is a constant $\flo>0$ such
that $\res(\st)\ge\flo$ for every state: no state is resolved to a point.
\end{theorem}

\begin{proof}[Proof sketch]
Resolving $\st$ to a point means completing the comparison of $\st$
against the whole of $\Man\setminus\set{\st}$. Suppose some state had zero
residue. Then a finite measurement completes that comparison, i.e.\
distinguishes $\st$ from all others at finite cost, which exhausts the
state's distinctions at a finite stage---contradicting non-exhaustibility
(\cref{ass:nonexhaust}). Hence every residue is positive. Over the
compact-by-resolution quotient of $\Man$ (states identified when
indistinguishable by \cref{ass:finite}) the positive residues are bounded
below by their infimum $\flo>0$.
\end{proof}

\begin{corollary}[Solved is a region]
\label{cor:region}
The finest determination of a state is a \emph{cell} of diameter $\flo$:
the set of states an ideal measurement cannot separate from the true one.
``The state'' is this cell, not its centre; a diagnosis is the
identification of the cell, and $\flo$ is the width within which the
question ``which state?'' has no finer answer.
\end{corollary}

\begin{principle}[The goal is the cell, not the point, and not a label]
\label{prin:goal}
Two consequences fix the whole method. (a) Because the target is a cell,
the diagnostic task is \emph{localization}---shrink the admissible region
to a single floor-width cell---not \emph{recognition against a catalogue}.
(b) Because the cell is defined intrinsically (by indistinguishability),
it needs no external library of named states to be identified; a diagnosis
can be complete without naming a disease.
\end{principle}

% =====================================================================
\section{Diagnosis as Localization}
\label{sec:localization}
% =====================================================================

We now make precise the sense in which a probe's trace \emph{locates} the
patient, and the condition under which a probe is diagnostically
sufficient.

\begin{definition}[State-resolved forward model]
\label{def:stateresolved}
A forward model of $\Sigma$ is \emph{state-resolved} if it predicts the
trace $\trace(\st_0,\probe)$ as a function of the initial state $\st_0$
over the whole manifold $\Man$, not merely at a distinguished reference
$\st^{\mathrm{ref}}$. Write $\Phi_\probe\colon\Man\to\mathcal{Y}^{[0,\infty)}$,
$\Phi_\probe(\st_0)=\trace(\st_0,\probe)$, for the predicted trace map of
probe $\probe$.
\end{definition}

\begin{remark}[Why ``reference-only'' is insufficient]
\label{rem:reference}
A model valid only at $\st^{\mathrm{ref}}$ predicts one trace; a measured
deviation from it certifies only \emph{that} the patient differs, not
\emph{how}. Inversion---reading the state \emph{off} the trace---requires
the trace to vary with the state, i.e.\ requires $\Phi_\probe$ as a map of
$\st_0$. This is the crux of the method and the sole substantive demand it
places on the underlying model.
\end{remark}

\begin{definition}[Localization map]
\label{def:loc}
Given a probe $\probe$ and a measured trace value $y\in\mathcal{Y}$ read at
time $t^\ast$, the \emph{localization} induced by $\probe$ is the preimage
\[
\loc_\probe(y)\;=\;\set{\st_0\in\Man :
\norm{\Rd(\Phi_\probe(\st_0)(t^\ast))-y}<\delta},
\]
the set of initial states whose predicted product distribution matches the
measurement to within the resolution $\delta$ (\cref{ass:finite}). Reading
the trace replaces the prior admissible region by $\loc_\probe(y)$.
\end{definition}

\begin{theorem}[Localization / diagnostic sufficiency]
\label{thm:localization}
A probe $\probe$ is \emph{diagnostically sufficient} for a region
$A\subseteq\Man$ iff the trace map restricted to $A$ is injective down to
the floor: for all $\st_0,\st_0'\in A$,
\[
\norm{\Rd(\Phi_\probe(\st_0)(t^\ast))-\Rd(\Phi_\probe(\st_0')(t^\ast))}<\delta
\;\Longrightarrow\;
\dist(\st_0,\st_0')<\flo .
\]
Under this condition, for every measurement $y$ the localization
$\loc_\probe(y)\cap A$ is contained in a single floor-width cell; the probe
pins the state to the resolution limit. Conversely, if the map is not
floor-injective on $A$, there exist distinct cells producing the same
reading, and no single reading of $\probe$ can separate them.
\end{theorem}

\begin{proof}[Proof sketch]
($\Leftarrow$) Floor-injectivity says equal readings force proximity below
$\flo$; hence the preimage of any reading has diameter $<\flo$ and lies in
one cell. ($\Rightarrow$) If two states at distance $\ge\flo$ share a
reading, their common reading's preimage meets two cells, so localization
fails to reach the floor. Sufficiency is exactly floor-injectivity of the
composed map $\st_0\mapsto\Rd(\Phi_\probe(\st_0)(t^\ast))$.
\end{proof}

\begin{corollary}[Probe design is an injectivity problem]
\label{cor:design}
Designing a diagnostic probe is choosing a probe (or a set) whose composed
trace map is floor-injective on the region of interest. Diagnosis is then
evaluation of the localization $\loc_\probe$ at the measured reading. No
disease library enters: the output is the cell, per \cref{prin:goal}.
\end{corollary}

% =====================================================================
\section{Cascading Probes and the Trajectory Certificate}
\label{sec:cascade}
% =====================================================================

A single reading of a single probe may fail to be floor-injective; more
information is needed. Rather than administer many independent probes, we
exploit that each metabolic product is itself a reactive species, so a
probe can be designed to generate its own successors---a cascade---whose
structure yields both richer localization and a strong non-forgeability
property.

\begin{definition}[Cascade probe]
\label{def:cascadeprobe}
A \emph{cascade probe} is a probe $\probe$ whose fate under $\fld$ passes
through a designed sequence of species $\probe=P_0\to P_1\to P_2\to\cdots$,
where each $P_{k+1}$ is a product of a reaction consuming $P_k$, and where
the reaction $P_k\to P_{k+1}$ is \emph{state-gated}: its rate (hence
whether $P_{k+1}$ forms in observable quantity within the read window)
depends on the state $\st$. The cascade is thus a state-dependent path
through the reaction graph, and the species $P_k$ present at read time
encode which transitions occurred.
\end{definition}

\begin{definition}[Trajectory; committed record]
\label{def:trajectory}
The \emph{traversed path} of a cascade in state $\st_0$ is the maximal
prefix $P_0\to\cdots\to P_{d(\st_0)}$ of state-gated transitions that
proceed to observable extent; $d(\st_0)$ is the \emph{depth}. Because each
transition consumes its predecessor and produces its successor, the
species population is a \emph{monotone committed record}: a later species'
presence is not undone by subsequent dynamics on the read timescale, so
the read distribution reports the path, not merely the instantaneous
state.
\end{definition}

\begin{theorem}[Trajectory certificate]
\label{thm:certificate}
Let $P_k$ be a species that forms in observable quantity only if every
transition $P_0\to\cdots\to P_k$ has occurred. Then observation of $P_k$ in
the trace \emph{entails} that the state $\st_0$ admits the whole path
$P_0\to\cdots\to P_k$: the set of states consistent with the reading is
contained in $\bigcap_{j\le k}\set{\st : \text{transition } j \text{ is
state-enabled at } \st}$. Consequently, a hypothesized state $\hat\st$ that
does not enable some transition $j\le k$ is \emph{refuted} by the presence
of $P_k$, and no such $\hat\st$ can reproduce the reading.
\end{theorem}

\begin{proof}
The forward implication is the existence condition of \cref{def:cascadeprobe}:
$P_k$ observable $\Rightarrow$ each prior transition proceeded
$\Rightarrow$ each prior transition was state-enabled at $\st_0$. The
consistent-state set is therefore the intersection of the enabling sets of
transitions $1,\dots,k$. Any $\hat\st$ outside transition $j$'s enabling set
predicts $P_j$, hence $P_k$, absent; observing $P_k$ contradicts $\hat\st$.
\end{proof}

\begin{principle}[Existence beats consistency]
\label{prin:existence}
\cref{thm:certificate} gives a stronger epistemic guarantee than agreement
of magnitudes. A consistency check asks whether a hypothesis \emph{matches}
observed values; it can be satisfied by a well-chosen falsehood. A cascade
reading asks whether the observed species \emph{could have formed}; a
falsehood that fails to enable a transition does not merely mismatch a
number---it predicts a species that is \emph{absent}, and absence cannot be
argued around. The terminal products are a constructive proof of their own
prerequisites.
\end{principle}

\begin{theorem}[Over-determination and rigidity]
\label{thm:rigidity}
Let a cascade (or a family of cascades) yield observed species
$S=\set{P_{k}}$. Define the \emph{surviving set}
$\Man_S=\bigcap_{P_k\in S}\set{\st:\text{path to }P_k\text{ is enabled at }\st}$.
As $S$ grows, $\Man_S$ is non-increasing. The state is
\emph{sufficiently localized} when $\Man_S$ has diameter $<\flo$: every
account of the state outside a single floor-width cell is refuted by the
absence it would predict of some observed species. This is a rigidity
condition---\emph{no alternative account survives}---not a threshold on any
scalar.
\end{theorem}

\begin{proof}
Each observed species adds a constraint $\set{\st:\text{its path
enabled}}$; the surviving set is their intersection, non-increasing in $S$
by monotonicity of intersection. When the intersection's diameter drops
below $\flo$ it lies in one cell (\cref{cor:region}); any $\st$ outside
that cell violates at least one constraint, i.e.\ predicts some observed
$P_k$ absent, and is refuted. Sufficiency is exactly diameter $<\flo$.
\end{proof}

\begin{corollary}[Single-measurement richness]
\label{cor:single}
A readout that returns the full species distribution---not one terminal
product but the population of all cascade products present---delivers the
entire constraint family $S$ in one measurement, hence the surviving set
$\Man_S$ in one shot. The instrument's report \emph{is} the committed
record of the traversed path.
\end{corollary}

\begin{remark}[Continuously working probe]
\label{rem:continuous}
Because the cascade is self-propagating and its record monotone, a single
administration installs a process that continues to advance and to record
after ingestion. The state is therefore readable at any later time by any
observer with the instrument: the probe does not deliver a value once and
expire; it remains a live, queryable source until the cascade completes.
``Continuous'' here is a property of the probe---always working---not of the
measurement schedule.
\end{remark}

% =====================================================================
\section{Therapy as Gap-Closure Without a Named Disease}
\label{sec:therapy}
% =====================================================================

We now define treatment purely in terms of the localized state and a
notion of coherence, with no reference to a disease label.

\begin{definition}[Coherent set; gap]
\label{def:coherent}
Let $\Man^\star\subseteq\Man$ be the \emph{coherent set}: the states in
which the pathway system is internally consistent (its conserved
quantities balance, its cyclic constraints close). For a localized state
$\st$, the \emph{gap} is the displacement to the nearest coherent state,
\[
\gap(\st)\;=\;\coh-\st,\qquad
\coh=\argmin_{z\in\Man^\star}\dist(\st,z).
\]
The gap is a measured vector, defined without naming what caused the
deviation.
\end{definition}

\begin{definition}[Intervention basis]
\label{def:basis}
An \emph{intervention basis} $\basis=\set{\intv_1,\dots,\intv_m}$ is a
collection of characterized compounds, each with a known state-resolved
effect: applying $\intv_i$ at state $\st$ produces a local displacement
$e_i(\st)\in\R^n$ (the compound's effect evaluated at the current state,
computed by the forward model). Compounds are indexed by \emph{properties}
(their displacement vectors), not by indication.
\end{definition}

\begin{problem}[Minimal gap-closing intervention]
\label{prob:closure}
Given a localized state $\st$ with gap $\gap(\st)$ and basis $\basis$,
choose nonnegative doses $c\in\Rp^{m}$ to close the gap with the fewest
compounds:
\[
\min_{c\in\Rp^{m}}\ \norm{c}_1
\quad\text{s.t.}\quad
\Bigl\lVert \gap(\st)-\textstyle\sum_{i}c_i\,e_i(\st)\Bigr\rVert\le\flo .
\]
The $\ell_1$ objective selects a sparse set of interventions; the
constraint requires the combined local displacement to close the gap to
within the floor.
\end{problem}

\begin{remark}[Repurposing is the default]
\label{rem:repurpose}
Because a compound enters \cref{prob:closure} through its displacement
$e_i(\st)$ and not through any indication, any compound whose displacement
has a component along the gap is admissible for any patient. The
distinction ``on-label / off-label'' does not arise: there are no labels,
only properties and gaps. The pharmacopeia is addressable by its effects.
\end{remark}

The local displacements $e_i(\st)$ are exact only at $\st$; applying the
chosen doses moves the state, after which the displacements change. Therapy
is therefore iterative.

\begin{definition}[Therapeutic map; contraction]
\label{def:contraction}
Let $T\colon\Man\to\Man$ send a state to the state resulting from
localizing it, solving \cref{prob:closure}, and applying the chosen
intervention. $T$ is a \emph{gap contraction} on a neighbourhood $U$ of
$\Man^\star$ if there is $\lambda\in[0,1)$ with
$\dist(T\st,\Man^\star)\le\lambda\,\dist(\st,\Man^\star)$ for $\st\in U$.
\end{definition}

\begin{theorem}[Convergence to coherence]
\label{thm:convergence}
If $T$ is a gap contraction on $U$ with modulus $\lambda$, then iterating
$T$ from any $\st_0\in U$ drives the state to the coherent set: the
sequence $\st_{k+1}=T\st_k$ satisfies
$\dist(\st_k,\Man^\star)\le\lambda^k\dist(\st_0,\Man^\star)\to 0$, and halts
when the distance falls below $\flo$. No global linearity of drug effects
is required---only that each measured, locally-composed step contracts the
gap.
\end{theorem}

\begin{proof}
Immediate by iterating the contraction inequality; the fixed point is in
$\overline{\Man^\star}$, and the floor provides the halting threshold below
which further contraction is unresolvable.
\end{proof}

\begin{principle}[Cure without naming]
\label{prin:cure}
\cref{thm:convergence} closes the gap---restores coherence---without ever
identifying a disease. The therapeutic objective is the measured residual,
and success is its reduction to the floor. Naming, if desired, is a
downstream description of which cell the state occupied; it is never a
prerequisite for treatment.
\end{principle}

% =====================================================================
\section{Diagnosis and Therapy Are One Operation}
\label{sec:unity}
% =====================================================================

The measurement primitive and the therapeutic primitive are the same
object read in two directions.

\begin{theorem}[Measurement--action identity]
\label{thm:unity}
Every characterized perturbation $\intv\in\mathcal{U}$ acts simultaneously
as a measurement and as an intervention:
\begin{enumerate}[label=\textup{(\roman*)},leftmargin=2.2em,itemsep=2pt]
\item as a \emph{measurement}, its trace $\trace(\st,\intv)$ localizes the
state (\cref{thm:localization});
\item as an \emph{intervention}, its displacement moves the state
(\cref{def:basis});
\end{enumerate}
and these are not two applications but one: applying $\intv$ yields a trace
(information) precisely by moving the state (action). A probe is a compound
read for its fate; a drug is a compound applied for its displacement; they
are the same primitive, and each does both.
\end{theorem}

\begin{proof}
A perturbation induces a trajectory (\cref{def:fate}); the trajectory's
readout is the measurement and the trajectory's endpoint is the action.
Both are functions of the single applied $\intv$. Hence the map
$\intv\mapsto(\text{trace},\ \text{displacement})$ has one input and the
two outputs are computed from the same induced trajectory.
\end{proof}

\begin{corollary}[Self-terminating therapeutic loop]
\label{cor:loop}
Because the therapeutic step is also a measurement, the contraction of
\cref{sec:therapy} needs no external monitor: each applied intervention
re-localizes the state as a side effect of acting. The loop
measure--act--re-measure is a single iterated perturbation, and it halts by
the rigidity criterion (\cref{thm:rigidity})---when the accumulated record
over-determines the state to within the floor---which is simultaneously the
diagnostic sufficiency condition and the therapeutic success condition.
\end{corollary}

\begin{principle}[Wellness as rigidity]
\label{prin:wellness}
Combining \cref{thm:rigidity,thm:convergence}: the terminal condition of
the loop is a state that is both \emph{coherent} (gap below floor) and
\emph{rigidly known} (no alternative account survives its own probes). This
double condition is the operative definition of wellness in the calculus:
a well system returns its probes at full resolution and admits no
gap-closing intervention beyond the floor.
\end{principle}

% =====================================================================
\section{The Runtime: Why It Must Be Derived, Not Fitted}
\label{sec:runtime}
% =====================================================================

Every construction above presupposes a forward model that predicts traces
across the manifold (state-resolution, \cref{def:stateresolved}). We now
state the condition this model must meet, as a theorem about
extrapolation, because it is the load-bearing requirement of the whole
method and the point at which most realizations fail.

\begin{definition}[Fitted vs.\ derived model]
\label{def:fitvsder}
A forward model is \emph{fitted} if its rate parameters are estimated to
minimize error against observations drawn from a training region
$A_{\mathrm{tr}}\subseteq\Man$. It is \emph{derived} if its rate laws and
their parameters are fixed by physical law (thermodynamic and kinetic
constants, conservation constraints) independently of any observation
region.
\end{definition}

\begin{theorem}[Extrapolation gap of fitted models]
\label{thm:extrap}
A fitted model carries no guarantee of accuracy outside
$A_{\mathrm{tr}}$: there exist state systems agreeing with the data on
$A_{\mathrm{tr}}$ to any tolerance yet disagreeing arbitrarily on
$\Man\setminus A_{\mathrm{tr}}$. Consequently a fitted model's predicted
trace map $\Phi_\probe$ is trustworthy only near $A_{\mathrm{tr}}$, whereas
a derived model, whose $\fld$ is fixed pointwise on $\Man$ by law, admits a
uniform error bound over $\Man$.
\end{theorem}

\begin{proof}[Proof sketch]
Standard non-identifiability: the data on $A_{\mathrm{tr}}$ constrain
$\fld$ only on $A_{\mathrm{tr}}$; any extension agreeing there is
admissible, and extensions may diverge off $A_{\mathrm{tr}}$. A derived
$\fld$ is not an extension chosen to fit but a function specified on all of
$\Man$ by its physical form, so its accuracy is governed by the fidelity of
the physical law, uniformly, not by proximity to data.
\end{proof}

\begin{principle}[Free experimentation requires a derived runtime]
\label{prin:derived}
The value of the method is that a user may experiment freely against the
model---searching the intervention space, testing cascades, exploring the
gap---at no risk to the patient. But free exploration \emph{is}
off-training-region evaluation by definition. By \cref{thm:extrap} a fitted
runtime becomes unreliable exactly where the user explores; only a derived
runtime supports risk-free \emph{in silico} search. The requirement is not
stylistic: the safety of the method reduces to the extrapolation validity
of the runtime.
\end{principle}

% =====================================================================
\section{The Language over the Runtime}
\label{sec:dsl}
% =====================================================================

Because the runtime is a runnable dynamical system and every diagnostic and
therapeutic act is a computation on it, the method presents as a
domain-specific language.

\begin{definition}[Core language]
\label{def:dsl}
The language $\dsl$ has three sorts---\emph{probes} (programs of reactions,
\cref{def:cascadeprobe}), \emph{states/regions} (subsets of $\Man$), and
\emph{intents} (target regions, typically $\Man^\star$)---and four
judgments:
\begin{enumerate}[label=\textup{(J\arabic*)},leftmargin=2.6em,itemsep=2pt]
\item \textsc{Run}: given a probe $\probe$ and a state $\st_0$, compute the
predicted trace $\Phi_\probe(\st_0)$ (integrate $\fld$).
\item \textsc{Localize}: given a probe $\probe$ and a measured reading $y$,
compute the localization $\loc_\probe(y)$ and, over a family of probes, the
surviving set $\Man_S$ (\cref{thm:rigidity}).
\item \textsc{Gap}: given a localized region, compute the gap $\gap$ to the
coherent set (\cref{def:coherent}).
\item \textsc{Close}: given a gap and a basis $\basis$, solve
\cref{prob:closure} for the minimal intervention and simulate the resulting
state (one step of $T$).
\end{enumerate}
A session is a program composing these judgments: \textsc{Run}/\textsc{Localize}
to a floor-width cell, \textsc{Gap}, \textsc{Close}, iterate.
\end{definition}

\begin{definition}[Soundness of the runtime]
\label{def:soundness}
The runtime is \emph{sound} for $\dsl$ if (i) \textsc{Run} agrees with the
true trace up to the uniform derived-model bound of \cref{thm:extrap};
(ii) \textsc{Localize} returns a superset of the true surviving set (no
false exclusion of the actual state); (iii) \textsc{Close} returns a
feasible intervention whose simulated displacement matches its physical
displacement up to the same bound. Soundness reduces to the derived-model
guarantee (\cref{prin:derived}).
\end{definition}

\begin{proposition}[Complexity]
\label{prop:complexity}
Let $n=\dim\Man$, $m=\abs\basis$, and let the cascade have depth $d$ and
the reading return $p$ species. Then: \textsc{Run} costs one numerical
integration of an $n$-dimensional system to the read time; \textsc{Localize}
costs $\mathcal{O}(p)$ constraint intersections against precomputed
enabling sets; \textsc{Gap} is a projection onto $\Man^\star$; and
\textsc{Close} is an $\ell_1$-regularized feasibility program in $m$
variables and $n$ constraints, solvable in polynomial time. The iterated
loop multiplies these by the number of contraction steps, which is
$\mathcal{O}(\log(1/\flo)/\log(1/\lambda))$ by \cref{thm:convergence}.
\end{proposition}

\begin{principle}[Relocation of the physician]
\label{prin:relocate}
Under $\dsl$, the mechanical work---predicting fates, inverting traces to
localizations, searching the intervention basis, simulating
compositions---is discharged by the runtime. What the language does
\emph{not} compute is the \emph{intent}: whether a given gap ought to be
closed, how far, and at what cost to the particular person. That judgment,
the selection of the target region among admissible coherent states, is the
irreducible human input. The physician is relocated from retrieval of
precedent to specification of intent over a derived, explorable state.
\end{principle}

% =====================================================================
\section{Constructive Validation}
\label{sec:validation}
% =====================================================================

Because every object is a dynamical system on a finite-dimensional manifold
and every quantity is an integration, a projection, or a sparse program,
each theorem is checkable by direct computation on explicit model
instances. The following protocol validates the calculus without appeal to
clinical data.

\begin{enumerate}[leftmargin=1.6em,itemsep=3pt]
\item \textbf{Floor} (\cref{thm:floor}). On a model $\Sigma$ with a fixed
readout resolution $\delta$, confirm that the induced state-discrimination
has a positive infimum $\flo$ and that no state is resolved below it.
\item \textbf{Localization} (\cref{thm:localization}). Sample state pairs;
confirm that a floor-injective probe never maps distinct cells to the same
reading within $\delta$, and exhibit a non-injective probe that does,
demonstrating insufficiency.
\item \textbf{Certificate} (\cref{thm:certificate}). Construct a state-gated
cascade; confirm that a species $P_k$ appears in the simulated trace iff
every prior transition is state-enabled, and that states failing a
transition produce $P_k$ absent.
\item \textbf{Rigidity} (\cref{thm:rigidity}). Grow the observed species set
and confirm the surviving set $\Man_S$ is non-increasing and reaches
diameter $<\flo$ exactly when no alternative cell survives.
\item \textbf{Gap-closure} (\cref{prob:closure}, \cref{thm:convergence}).
On states off the coherent set, solve the sparse program and iterate $T$;
confirm geometric contraction of $\dist(\cdot,\Man^\star)$ and sparsity of
the selected intervention set.
\item \textbf{Unity} (\cref{thm:unity}). Confirm that a single simulated
perturbation yields both a localizing trace and a state displacement from
one integration.
\item \textbf{Derived vs.\ fitted} (\cref{thm:extrap}). Fit a surrogate model
on a training region, then evaluate both surrogate and derived model off
that region; confirm the surrogate's trace error grows off-region while the
derived model's stays within its uniform bound---quantifying why free
exploration requires the derived runtime.
\end{enumerate}

Each check is a finite, deterministic computation whose predicted and
measured outcomes agree exactly for identities and within the stated bound
for inequalities. The suite validates the calculus on explicit model
instances; its fidelity to a real organism is inherited entirely from the
fidelity of the derived forward model, which is the separate, physical
question the runtime must answer.

% =====================================================================
\section{Summary}
\label{sec:summary}
% =====================================================================

We have described a calculus of medicine built on one primitive---a
characterized perturbation of metabolism---and one model requirement---that
the forward model predicting the perturbation's fate be state-resolved and
derived rather than fitted. From a positive resolution floor we obtained
that a state is a region, so diagnosis is localization to a floor-width cell
rather than recognition against a disease catalogue. A probe's trace
localizes the state when its trace map is floor-injective; a cascading
probe strengthens this to a trajectory certificate, since its later
products exist only if the earlier path occurred, so that a false account
cannot forge the reading and the correct stopping rule is rigidity---halt
when no alternative account survives. Therapy is the closure of the measured
gap to the coherent set by a sparse combination of characterized
interventions selected by property, executed as a contraction to a fixed
point, with no disease ever named. Diagnosis and therapy coincide, because
every characterized perturbation both measures and moves. Finally, the
runnable forward model makes the whole method a domain-specific language: a
user specifies probes and intents, the runtime returns traces,
localizations, gaps, and minimal interventions, and---because the runtime
is derived---the user may search the intervention space freely and safely
\emph{in silico}, applying to the patient only the intervention the search
selects. The physician's role moves from precedent-retrieval to the
specification of intent over a derived, explorable state.

\end{document}
